\documentclass[10pt,a4paper]{amsart}
\usepackage[margin=19mm]{geometry}
\usepackage{amsmath,amssymb,mathtools}
\usepackage[scr=rsfs,cal=euler]{mathalfa}
\usepackage{xfrac}
\usepackage[unicode,hidelinks,bookmarks=false]{hyperref}
\newcommand{\ie}{i.e.\ }
\newcommand{\cf}{cf.\ }
\newcommand{\resp}{respectively\ }
\newcommand{\Subsection}{\subsection}
\providecommand{\nobreakdash}{\nobreak\hbox{-}}
\setlength{\emergencystretch}{2em}
\multlinegap0pt
\usepackage{amssymb}
\usepackage{algorithm}\usepackage{algpseudocode}
\newtheorem{thm}{Theorem}
\def\cal{\mathcal}
\title{There are infinitely many Hilbert cubes of dimension 3 in the set of squares}
\author{Andrew Bremner, Christian Elsholtz, Maciej Ulas}
\date{Source: arXiv:2604.05459v1 (CC BY 4.0)}
\begin{document}
\maketitle

\noindent\emph{Source and licence.} There are infinitely many Hilbert cubes of dimension 3 in the set of squares. Version arXiv:2604.05459v1 (CC BY 4.0), distributed by arXiv under the Creative Commons Attribution 4.0 International licence, http://creativecommons.org/licenses/by/4.0/, as recorded in the arXiv licence metadata for this exact version and shown on the abstract page https://arxiv.org/abs/2604.05459v1. Full licence notice: \texttt{licenses/arxiv-2604.05459-CC-BY-4.0.txt}.\par\smallskip

\noindent\textit{Editorial scope.} Counted unit: the article's thm thm (source0.tex lines 1037--1039). The article's complete original proof of it is delivered with it (source0.tex lines 1040--1118, one \texttt{proof} environment of 79 lines). No mathematics is written, repaired or re-derived: the statement and the proof are the article's own lines, in the article's own notation, and no retained line is substituted or re-flowed.  No further source-local context is needed: every cross-reference of every retained line resolves inside the two delivered spans.  Nothing was omitted from any retained span, so the package carries no omission record.  No retained line contains a citation, so the wrapper carries no bibliography and no bibliography entry.  No retained line contains a figure environment, an image-inclusion command or drawing code, so no image is delivered and the package declares no asset dependency.  Editorial formatting only, in addition: an AMS \texttt{amsart} preamble that restates the article's own declarations and macros used by the retained lines, namely the environments \texttt{thm} on separate counters, together with the standard packages those lines need.  The retained environments print in this document as Theorem 1 in order of appearance, whereas the article prints the same items with the numbering of its own full document.  The complete native source of the article as distributed in the arXiv e-print for this exact version is bundled unchanged as \texttt{sources/197959/source0.tex}.
\begin{thm}
    There are infinitely many pairs $H(a_{0}; a_{1}, a_{2}, a_{3}), H(A_0; A_1, A_2, A_3)\subset \cal{S}$ of Hilbert subes such that $a_{i}=A_{i}$ for $i=0, 1, 2$.
\end{thm}

\begin{proof} 
The defining system is
\begin{center}
\begin{tabular}{rrrr}
	$a_0=p^2$ & $a_0+a_1=q^2$ & $a_0+a_2=r^2$ & $a_0+a_1+a_2+a_3=s^2$ \\
	$a_0+a_3=P_1^2$ & $a_0+a_1+a_3=Q_1^2$ & $a_0+a_2+a_3=R_1^2$ & $a_0+a_1+a_2+a_3=S_1^2$ \\
	$a_0+A_3=P_2^2$ & $a_0+a_1+A_3=Q_2^2$ & $a_0+a_2+A_3=R_2^2$ & $a_0+a_1+a_2+A_3=S_2^2$.
\end{tabular}
\end{center}
On eliminating $a_0,a_1,a_2,a_3,A_3$, this is equivalent to the system
\[ p^2+s^2=q^2+r^2, \qquad P_1^2+S_1^2=Q_1^2+R_1^2, \qquad P_2^2+S_2^2=Q_2^2+R_2^2, \]
\[ p^2-q^2 = P_1^2-Q_1^2 = P_2^2-Q_2^2, \qquad p^2-r^2 = P_1^2-R_1^2 = P_2^2-R_2^2. \]
To satisfy the first row, set
\begin{align*}
	(p,q,r,s) = & (\alpha_0 \delta_0-\beta_0 \gamma_0, \;\; \alpha_0 \delta_0+\beta_0 \gamma_0, \;\; \alpha_0 \gamma_0-\beta_0 \delta_0, \;\; \alpha_0 \gamma_0+\beta_0 \delta_0), \\
	(P_1,Q_1,R_1,S_1) = & (\alpha_1 \delta_1-\beta_1 \gamma_1, \;\; \alpha_1
 \delta_1+\beta_1 \gamma_1, \;\; \alpha_1 \gamma_1-\beta_1 \delta_1, \;\; \alpha_1 \gamma_1+\beta_1 \delta_1), \\
	(P_2,Q_2,R_2,S_2) = & (\alpha_2 \delta_2-\beta_2 \gamma_2, \;\; \alpha_2
 \delta_2+\beta_2 \gamma_2, \;\; \alpha_2 \gamma_2-\beta_2 \delta_2, \;\; \alpha_2 \gamma_2+\beta_2 \delta_2).
\end{align*}
The second row then delivers
\[ \alpha_0 \beta_0 \gamma_0 \delta_0 = \alpha_1 \beta_1 \gamma_1 \delta_1 = \alpha_2 \beta_2 \gamma_2 \delta_2, \]
\[ (\alpha_0^2-\beta_0^2)(\gamma_0^2-\delta_0^2) = (\alpha_1^2-\beta_1^2)(\gamma
_1^2-\delta_1^2) = (\alpha_2^2-\beta_2^2)(\gamma_2^2-\delta_2^2). \] 
If we set
\[ (\alpha_i,\beta_i)=\left( \frac{\gamma_i+\delta_i}{2},\;\frac{\gamma_i-\delta_i}{2} \right), \quad i=0,1,2, \]
then we obtain the simple system
\[ \gamma_0 \delta_0(\gamma_0^2-\delta_0^2) = \gamma_1 \delta_1(\gamma_1^2-\delta_1^2) = \gamma_2 \delta_2(\gamma_2^2-\delta_2^2), \]
requiring three Pythagorean triangles with equal area. Such are provided by
\begin{align*}
	(\gamma_0,\delta_0) = & (u^2+u v+v^2, u^2-v^2), \\
	(\gamma_1,\delta_1) = & (u^2+u v+v^2, 2u v+v^2), \\
	(\gamma_2,\delta_2) = & (u^2+2 u v, u^2+u v+v^2),
\end{align*}
with 
\begin{align*}
	(\alpha_0,\beta_0) = & (u(2u+v)/2, v(u+2v)/2), \\
	(\alpha_1,\beta_1) = & ((u+v)(u+2v)/2, u(u-v)/2) \\
	(\alpha_2,\beta_2) = & ((u+v)(2u+v)/2, (u-v)v/2).
\end{align*}
Then
\begin{align*}
	(p,q,r,s) = & ((2u^4 - 5u^2v^2 - 4uv^3 - 2v^4)/2, \\
	& (2u^4 + 2u^3v + u^2v^2 + 2uv^3 + 2v^4)/2, \\
	& (2u^4 + 2u^3v + u^2v^2 + 2uv^3 + 2v^4)/2, \\
	& (2u^4 + 4u^3v + 5u^2v^2 - 2v^4)/2);
\end{align*}
% {(2*u^4 - 5*u^2*v^2 - 4*u*v^3 - 2*v^4)/2, (2*u^4 + 2*u^3*v + u^2*v^2 + 2*u*v^3 + 2*v^4)/2, (2*u^4 + 2*u^3*v + u^2*v^2 + 2*u*v^3 + 2*v^4)/2, (2*u^4 + 4*u^3*v + 5*u^2*v^2 - 2*v^4)/2}
\begin{align*}
	(P_1,Q_1,R_1,S_1) = & ((-u^4 + 2u^3v + 7u^2v^2 + 8uv^3 + 2v^4)/2, \\
	& (u^4 + 2u^3v + 7u^2v^2 + 6uv^3 + 2v^4)/2,  \\
	& (u^4 + 2u^3v + 7u^2v^2 + 6uv^3 + 2v^4)/2, \\
	& (u^4 + 6u^3v + 5u^2v^2 + 4uv^3 + 2v^4)/2)
\end{align*}
% {(-u^4 + 2*u^3*v + 7*u^2*v^2 + 8*u*v^3 + 2*v^4)/2, (u^4 + 2*u^3*v + 7*u^2*v^2 + 6*u*v^3 + 2*v^4)/2, (u^4 + 2*u^3*v + 7*u^2*v^2 + 6*u*v^3 + 2*v^4)/2, (u^4 + 6*u^3*v + 5*u^2*v^2 + 4*u*v^3 + 2*v^4)/2}
\begin{align*}
	(P_2,Q_2,R_2,S_2) = & ((2u^4 + 4u^3v + 5u^2v^2 + 6uv^3 + v^4)/2, \\
	& (2u^4 + 6u^3v + 7u^2v^2 + 2uv^3 + v^4)/2, \\
	& (2u^4 + 6u^3v + 7u^2v^2 + 2uv^3 + v^4)/2, \\
	& (2u^4 + 8u^3v + 7u^2v^2 + 2uv^3 - v^4)/2).
\end{align*}
% {(2*u^4 + 4*u^3*v + 5*u^2*v^2 + 6*u*v^3 + v^4)/2, (2*u^4 + 6*u^3*v + 7*u^2*v^2 + 2*u*v^3 + v^4)/2, (2*u^4 + 6*u^3*v + 7*u^2*v^2 + 2*u*v^3 + v^4)/2, (2*u^4 + 8*u^3*v + 7*u^2*v^2 + 2*u*v^3 - v^4)/2}
Finally, scaling by a factor 4, 
\begin{align*}
	(a_0,a_1,a_2,a_3) = &  ( (2u^4 - 5u^2v^2 - 4uv^3 - 2v^4)^2, \\
	& 4u(u - v)v(u + v)(2u + v)(u + 2v)(u^2 + uv + v^2), \\
	& 4u(u - v)v(u + v)(2u + v)(u + 2v)(u^2 + uv + v^2), \\
	& -u(u + 2v)(u^2 - 2uv - 2v^2)(u^2 + 2v^2)(3u^2 + 4uv + 2v^2) );
\end{align*}
% {(2*u^4 - 5*u^2*v^2 - 4*u*v^3 - 2*v^4)^2/4, u*(u - v)*v*(u + v)*(2*u + v)* (u+ 2*v)*(u^2 + u*v + v^2), u*(u - v)*v*(u + v)*(2*u + v)*(u + 2*v)* (u^2 + u*v + v^2), -1/4*(u*(u + 2*v)*(u^2 - 2*u*v - 2*v^2)*(u^2 + 2*v^2)* (3*u^2 + 4*u*v + 2*v^2))}
\begin{align*}
	(A_0,A_1,A_2,A_3) = & ((2u^4 - 5u^2v^2 - 4uv^3 - 2v^4)^2, \\
	& 4u(u - v)v(u + v)(2u + v)(u + 2v)(u^2 + uv + v^2), \\
	& 4u(u - v)v(u + v)(2u + v)(u + 2v)(u^2 + uv + v^2), \\
	& v(2u + v)(2u^2 + 2uv - v^2)(2u^2 + v^2)(2u^2 + 4uv + 3v^2)),
\end{align*}
% {(2*u^4 - 5*u^2*v^2 - 4*u*v^3 - 2*v^4)^2/4, u*(u - v)*v*(u + v)*(2*u + v)* (u+ 2*v)*(u^2 + u*v + v^2), u*(u - v)*v*(u + v)*(2*u + v)*(u + 2*v)* (u^2 + u*v +v^2), (v*(2*u + v)*(2*u^2 + 2*u*v - v^2)*(2*u^2 + v^2)* (2*u^2 + 4*u*v + 3*v^2))/4}
where $a_0=A_0$, $a_1=A_1$, $a_2=A_2$.
\end{proof}

\end{document}
